
\documentclass[11pt,a4paper]{article}

\usepackage[margin=1in]{geometry}
\usepackage{amsmath,amssymb,amsthm}
\usepackage{enumitem}

\setlength{\parindent}{0pt}
\setlength{\parskip}{5pt}

\newcommand{\R}{\mathbb{R}}
\newcommand{\N}{\mathbb{N}}
\newcommand{\Rw}{\mathbb{R}^{\omega}}
\newcommand{\Rl}{\mathbb{R}_{\ell}}
\newcommand{\Bdd}{\mathcal{B}}
\newcommand{\Unb}{\mathcal{U}}

\newtheorem*{problem}{Problem}

\begin{document}

\begin{center}
    {\Large\textbf{Solutions to Problems 5--8}}
\end{center}

%---------------- PROBLEM 5 ----------------

\begin{problem}
Let $(X,d)$ be a metric space and let
$S\subseteq X$ be non-empty. Define
\[
d(x,S)=\inf_{s\in S}d(x,s).
\]
\begin{enumerate}[label=(\arabic*)]
    \item Show that
    \[
    |d(x,S)-d(y,S)|\leq d(x,y).
    \]
    \item Prove that $f_S:X\to\R$, defined by
    $f_S(x)=d(x,S)$, is continuous.
\end{enumerate}
\end{problem}

\textbf{Proof.}

(1) By the triangle inequality, for every $s\in S$,
\[
d(x,s)\leq d(x,y)+d(y,s).
\]
Taking the infimum over $s\in S$, we get
\[
d(x,S)\leq d(x,y)+d(y,S).
\]
Thus,
\[
d(x,S)-d(y,S)\leq d(x,y).
\]
Interchanging $x$ and $y$ gives
\[
d(y,S)-d(x,S)\leq d(x,y).
\]
Therefore,
\[
\boxed{|d(x,S)-d(y,S)|\leq d(x,y).}
\]

(2) Let $x_0\in X$ and $\varepsilon>0$.
Choose $\delta=\varepsilon$.

If $d(x,x_0)<\delta$, then by (1),
\[
\begin{aligned}
|f_S(x)-f_S(x_0)|
&=|d(x,S)-d(x_0,S)|\\
&\leq d(x,x_0)\\
&<\varepsilon.
\end{aligned}
\]
Hence, $f_S$ is continuous on $X$.
\hfill $\square$

\vspace{0.5cm}

%---------------- PROBLEM 6 ----------------

\begin{problem}
Let $\{A_n\}_{n\in\N}$ be a sequence of connected
subspaces of $X$ such that
\[
A_n\cap A_{n+1}\neq\emptyset
\quad\text{for all }n\in\N.
\]
Prove that $\displaystyle\bigcup_{n=1}^{\infty}A_n$
is connected.
\end{problem}

\textbf{Proof.}

Let
\[
A=\bigcup_{n=1}^{\infty}A_n.
\]
Suppose, for contradiction, that $A$ is disconnected.
Then
\[
A=U\cup V,
\]
where $U,V$ are non-empty, disjoint and open in $A$.

Since each $A_n$ is connected and
\[
A_n=(A_n\cap U)\cup(A_n\cap V),
\]
we must have
\[
A_n\subseteq U
\quad\text{or}\quad
A_n\subseteq V.
\]

Without loss of generality, let $A_1\subseteq U$.

Since $A_1\cap A_2\neq\emptyset$, there exists
$p\in A_1\cap A_2$. Thus $p\in U$, so
$A_2\not\subseteq V$. Consequently,
$A_2\subseteq U$.

Repeating this argument inductively, we obtain
\[
A_n\subseteq U
\quad\text{for every }n\in\N.
\]

Therefore,
\[
A=\bigcup_{n=1}^{\infty}A_n\subseteq U,
\]
contradicting $V\neq\emptyset$.

Hence,
\[
\boxed{\bigcup_{n=1}^{\infty}A_n
\text{ is connected.}}
\]
\hfill $\square$

\vspace{0.5cm}

%---------------- PROBLEM 7 ----------------

\begin{problem}
Determine whether or not $\Rw$ is connected
under the uniform topology.
\end{problem}

\textbf{Proof.}

The uniform metric is
\[
\bar{\rho}(x,y)
=\sup_n\min\{|x_n-y_n|,1\}.
\]

Define
\[
\Bdd=
\left\{x\in\Rw:\sup_n|x_n|<\infty\right\}
\]
and
\[
\Unb=\Rw\setminus\Bdd.
\]

We show that both sets are open.

First, observe that
\[
\bar{\rho}(x,y)<1
\implies |x_n-y_n|<1
\quad\forall n.
\]

\textbf{Claim 1: $\Bdd$ is open.}

Let $x\in\Bdd$. Then $|x_n|\leq M$ for
some $M>0$ and all $n$.

If $\bar{\rho}(x,y)<1$, then
\[
|y_n|\leq |x_n|+|y_n-x_n|<M+1.
\]
Hence $y\in\Bdd$, proving that $\Bdd$ is open.

\textbf{Claim 2: $\Unb$ is open.}

Let $x\in\Unb$ and suppose
$\bar{\rho}(x,y)<1$.

If $y$ were bounded, say $|y_n|\leq M$,
then
\[
|x_n|\leq |y_n|+|x_n-y_n|<M+1.
\]
This contradicts the fact that $x$ is unbounded.

Thus $y\in\Unb$, proving that $\Unb$ is open.

Finally, $\Bdd$ and $\Unb$ are disjoint,
non-empty, and
\[
\Bdd\cup\Unb=\Rw.
\]

Therefore, they form a separation of $\Rw$.
Hence,
\[
\boxed{\Rw\text{ is not connected.}}
\]
\hfill $\square$

\vspace{0.5cm}

%---------------- PROBLEM 8 ----------------

\begin{problem}
Determine whether or not the lower limit topology
$\Rl$ on $\R$ is connected.
\end{problem}

\textbf{Proof.}

The lower limit topology has the basis
\[
\mathcal{B}=\{[a,b):a<b\}.
\]

Consider
\[
U=(-\infty,0),\qquad V=[0,\infty).
\]

We show that both sets are open.

If $x\in U$, then
\[
x\in[x,0)\subseteq U.
\]
Since $[x,0)$ is a basis element, $U$ is open.

Similarly, if $x\in V$, then
\[
x\in[x,x+1)\subseteq V.
\]
Thus, $V$ is also open.

Moreover,
\[
U\cap V=\emptyset,
\qquad U\cup V=\R.
\]

Both sets are non-empty. Therefore, $(U,V)$
forms a separation of $\Rl$.

Hence,
\[
\boxed{\Rl\text{ is not connected.}}
\]
\hfill $\square$

\end{document}
